\documentclass[10pt,a4paper]{amsart}
\usepackage[margin=19mm]{geometry}
\usepackage{amsmath,amssymb,mathtools}
\usepackage[scr=rsfs,cal=euler]{mathalfa}
\usepackage{xfrac}
\usepackage[unicode,hidelinks,bookmarks=false]{hyperref}
\newcommand{\ie}{i.e.\ }
\newcommand{\cf}{cf.\ }
\newcommand{\resp}{respectively\ }
\newcommand{\Subsection}{\subsection}
\providecommand{\nobreakdash}{\nobreak\hbox{-}}
\setlength{\emergencystretch}{2em}
\multlinegap0pt
\usepackage{amssymb}
\usepackage{mathtools}
\usepackage[shortlabels]{enumitem}
\newtheorem{lemma}{Lemma}
\newcommand{\C}{\mathbb{C}}
\newcommand{\h}{\mathfrak{h}}
\newcommand{\n}{\mathfrak{n}}
\title{Equivariant compactifications of a unipotent group by a smooth projective horospherical variety of Picard number one}
\author{Hyukmoon Choi}
\date{Source: arXiv:2609.08072v1 (CC BY 4.0)}
\begin{document}
\maketitle

\noindent\emph{Source and licence.} Equivariant compactifications of a unipotent group by a smooth projective horospherical variety of Picard number one. Version arXiv:2609.08072v1 (CC BY 4.0), distributed by arXiv under the Creative Commons Attribution 4.0 International licence, http://creativecommons.org/licenses/by/4.0/, as recorded in the arXiv licence metadata for this exact version and shown on the abstract page https://arxiv.org/abs/2609.08072v1. Full licence notice: \texttt{licenses/arxiv-2609.08072-CC-BY-4.0.txt}.\par\smallskip

\noindent\textit{Editorial scope.} Counted unit: the article's lemma unipotent (source0.tex lines 187--192). The article's complete original proof of it is delivered with it (source0.tex lines 194--226, one \texttt{proof} environment of 33 lines). No mathematics is written, repaired or re-derived: the statement and the proof are the article's own lines, in the article's own notation, and no retained line is substituted or re-flowed.  No further source-local context is needed: every cross-reference of every retained line resolves inside the two delivered spans.  Nothing was omitted from any retained span, so the package carries no omission record.  No retained line contains a citation, so the wrapper carries no bibliography and no bibliography entry.  No retained line contains a figure environment, an image-inclusion command or drawing code, so no image is delivered and the package declares no asset dependency.  Editorial formatting only, in addition: an AMS \texttt{amsart} preamble that restates the article's own declarations and macros used by the retained lines, namely the environments \texttt{lemma} on separate counters, together with the standard packages those lines need.  The retained environments print in this document as Lemma 1 in order of appearance, whereas the article prints the same items with the numbering of its own full document.  The complete native source of the article as distributed in the arXiv e-print for this exact version is bundled unchanged as \texttt{sources/209694/source0.tex}.
\begin{lemma}
\label{unipotent}
In the above notation, $N$ is a subgroup of $H^-$ satisfying
\[N \cong H^- /(H \cap H^-)\]
as an algebraic variety. 
\end{lemma}

\begin{proof}
We first show
\[H_u^- = (H_u \cap H_u^-) \cdot N \text{, and } N \cong H_u^- /(H_u \cap H_u^-)\]
and then show that
\[H_u^- /(H_u \cap H_u^-) \cong H^- /(H \cap H^-)\]
as an algebraic variety.

Since the type of $L$ is one of $B_n$, $C_n$, $F_4$, and $G_2$, $w_0(\omega) = -\omega$ for any weight $\omega$ of $L$. Hence, we have
\begin{align*}
\h^- &= (\bigoplus_{k\leq0} \mathfrak{l}_k \oplus \C) \rhd \bigoplus_{k = -\mu}^{\mu-1}V_k, \\
\h \cap \h^-  &= (\mathfrak{l}_0 \oplus \C) \rhd \bigoplus_{k = -\mu+1}^{\mu-1}V_k,\\
\h_u &= \bigoplus_{k > 0} \mathfrak{l}_k \rhd \bigoplus_{k = -\mu + 1}^{\mu}V_k,\\
\h_u^- &= \bigoplus_{k < 0} \mathfrak{l}_k \rhd \bigoplus_{k = -\mu}^{\mu-1}V_k.
\end{align*}

The Lie algebra of $H_u \cap H_u^-$ is equal to 
\[\h_u \cap \h_u^- = \bigoplus_{k=-\mu+1}^{\mu-1} V_{k}.\]
On the other hand, $N$ is a nilpotent group with Lie algebra
\[\n = (\bigoplus_{k < 0} \mathfrak{l}_k) \oplus V_{-\mu}.\]
Hence, we have 
\[\h_u^- = \n + (\h_u \cap \h_u^-) \text{ and } \n = \h_u^- / (\h_u \cap \h_u^-)\]
as vector spaces, which implies that
\[H_u^- = (H_u \cap H_u^-) \cdot N \text{ and } N\cong H_u^-/(H_u \cap H_u^-)\]
as algebraic varieties.

The explicit descriptions of the Lie algebras above show that
\[H_u \cap H_u^- = (H\cap H^-) \cap H_u^-\]
and 
\[(H \cap H^-) \cdot H_u^- = H^-.\]
Hence, $H_u^-$ acts on $H^-/(H \cap H^-)$ transitively and the isotropy group at the identity coset is $H_u \cap H_u^-$. Consequently, 
\[H_u^- /(H_u \cap H_u^-) \cong H^- /(H \cap H^-)\]
as algebraic varieties.
\end{proof}

\end{document}
